\section{Trimming: DST vs.\ trimmed DST vs.\ ntuple}
\label{sec:size}

\begin{table}[H]
\centering\small
\input{tables/tab_dstsize.tex}
\caption{Compressed size per event of each node of the production DST and of the trimmed DST
(test segment, \nEvents{} events; ROOT \fcmd{GetZipBytes} of the branches of tree \fcmd{T}, run tree
not included). The trimmed DST is written by \fcmd{Fun4AllDstOutputManager} with its default
compression, which makes \fcmd{SvtxTrackMap} slightly larger than in the input. Source:
\fpath{logs/dataset_software.txt}.}
\label{tab:dstsize}
\end{table}

Clusters (\fcmd{TRKR_CLUSTER_SEED}) are \pctClus\% and the three seed containers \pctSeeds\% of
the DST; the track map is \pctTrkMap\% (\trkMapkBPerEv\,kB/event) and the vertex map
\pctVtxMap\%. File sizes per event:

\begin{center}\small
\begin{tabular}{lrrr}
\toprule
file & size (MB) & kB/event & reduction\\
\midrule
production DST & \sizeDSTMB & \sizeDSTkB & 1\\
trimmed DST (track + vertex maps, GL1, header) & \sizeTrimMB & \sizeTrimkB & \ratioTrim\\
TrackQA ntuple & \sizeNtMB & \sizeNtkB & \ratioNt\\
\bottomrule
\end{tabular}
\end{center}

\choice{Trimmed DST = \fcmd{Sync}, \fcmd{EventHeader}, \fcmd{GL1RAWHIT}, \fcmd{SvtxTrackMap},
\fcmd{SvtxVertexMap}. It keeps the full \fcmd{SvtxTrack} objects (states, covariances, seed ids),
so modules that need only tracks and vertices can run on it; for \fcmd{TrackQATree} this was
verified: re-running \fcmd{TrackQATree} on it reproduces the
ntuple exactly (check \fcmd{chkTrim}). The file name keeps the production pattern
\fpath{DST_TRKR_TRACKS_TRIM-<run>-<seg>.root} so that \fcmd{Fun4AllUtils::GetRunSegment} works on it.}

\caveat{Anything that needs clusters or seeds (dE/dx, residuals, refits, KFParticle detailed
tracking or dE/dx branches) does not work on the trimmed DST. Running KFParticle on it was not
tested.}
